\section{The last bounded-shape C2 corner: a weighted radial root}
\label{sec:c2-radial-corner}

It remains to treat the corner in which $h,r,g,j$ are all positive
infinitesimals.  The radial coefficient then tends to zero, but the
positive source identity gives exactly the weight needed to normalize it.

\begin{theorem}[The all-small radial corner]
\label{thm:c2-radial-corner}
Use the unchanged source \eqref{eq:c2u-source}.  Suppose that $c>0$ and
$h=b-a^2t,r,g,j>0$ are infinitesimal, while $a,1-r,t$ are positive
bounded units.  For every independent infinitesimal or zero testing pair
$u,v$ and every real or complex projective state, the full parent is a
finite sum of squares in $K_c\otimes_{\mathbb R}K_t$.
\end{theorem}

\subsection{The small source ideal controls all coefficients}
Write $z=c^2$ and $D_c=1+c^2$.  The source identities give
\begin{equation}
b-1=-r+g-zh,\qquad t-1=zr-j,\qquad
a-1=\frac{g+j-D_c(r+h)}{t(a+1)}.
\label{eq:c2r-source-ideal}
\end{equation}
Thus $a,b,t$ all have residue $1$.  Let $I=(h,r,g,j)$ in the ring of
bounded source coefficients used below.  By
\eqref{eq:c2y-N-positive}, $N>0$ and $N/(h+r+g+j)$ has residue $1$.
In particular
\begin{equation}
h/N,\quad r/N,\quad g/N,\quad j/N
\text{ are bounded},\qquad
4\alpha/(c^4N)\text{ has residue }1.
\label{eq:c2r-weight-bounds}
\end{equation}

All rational coefficient functions in the retained remainder have
denominators with nonzero residue: the possibly small denominators
$r,g,j,h$ have been absorbed into $U_1,U_2,U_5,U_4$.  Specialize a
bounded rational coefficient by putting $h=r=g=j=0$ and $a=b=t=1$,
leaving $c$ untouched.  Its difference from this specialization lies in
$I$, with bounded coquotients.  Indeed finite polynomial subtraction
factors the differences by $r,g,j,a-1$, and
\eqref{eq:c2r-source-ideal} replaces $a-1$ by bounded multiples of the
four generators.  Subtraction of fractions only uses unit denominators.
The monic polynomial divisions in the preceding proofs preserve this
property.  Pure-testing denominators $D_u,D_v$ do so as well.
Consequently every coefficient in $I$, divided by the pure coefficient
$N$, remains bounded.  This observation involves no comparison with
the independent testing infinitesimals.

\subsection{An explicit additional zero correction}
Use the four retained packets and six-root Gram from
Section~\ref{sec:c2-y-current}, including its zero correction
\eqref{eq:c2y-zero-correction}.  That construction required only that
$e_3$ be bounded; it still gives an exact bounded Gram here.  Two of its
current coefficients simplify exactly to
\begin{equation}
Q(u^2)=h-a^2ru^2,\qquad Q(u^2)+a^2D_ur=h+a^2r.
\label{eq:c2r-QJ}
\end{equation}
Both belong to $I$.  The coefficient of $T$ in the exact relation
$\zeta$ is therefore $u(h+a^2r)$.

We record the one base calculation needed for the new correction.
On the scalar normal graph, after specialization modulo $I$, put
\[
Y_*=uW-v,\qquad G_*=D_uvW-D_vu.
\]
The entire bounded scalar remainder becomes
\begin{equation}
f_*=2uG_*+\frac{2}{D_c}D_uD_vuY_*
 +2\{D_u(vY_*-1)^2+D_vu^2-D_uD_v\}.
\label{eq:c2r-base-polynomial}
\end{equation}
Here $U_2=1/4$, $U_4=1/(4D_c)$; all terms containing a generator of
$I$, $e_3$, or $\alpha$ vanish.  Direct differentiation gives
\begin{equation}
\begin{gathered}
f_*(u,u,1)=0,\\
(f_*)_W(u,u,1)=\frac{2u^2(u^2-c^2)D_u}{D_c},\qquad
(f_*)_v(u,u,1)=-\frac{2u(u^2-c^2)D_u}{D_c}.
\end{gathered}
\label{eq:c2r-base-derivatives}
\end{equation}

Use the division order of the preceding proof: first restrict to $v=au$,
then expand at $W_c$, and retain the finite second-order remainder in
$v-au$ as an $R^2$ coefficient.  The first two derivatives in
\eqref{eq:c2r-base-derivatives} imply that the $ST$ coefficient is
$2uD_u/D_c$ and the direct $RT$ coefficient is $2uD_u/D_c$.
Modulo $I$, elimination of $S$ reads
\[
S=\frac{H_2+(1-u^2v^2)H_1+uvH_5
 +(D_uW-u(v+u))R}{D_u}.
\]
Hence the full scalar $RT$ coefficient after elimination is exactly
\begin{equation}
\frac{2u}{D_c}(D_uW+1-uv).
\label{eq:c2r-base-RT}
\end{equation}
The earlier zero correction adds no $RT$ term at this specialization,
because \eqref{eq:c2r-QJ} vanishes there.  Meanwhile the coefficient of
$R$ in $\zeta$ specializes to $-(D_uW+1-uv)$.  Add one more exact zero:
\begin{equation}
\frac{2u}{D_c}\,\zeta T=0.
\label{eq:c2r-final-correction}
\end{equation}
By \eqref{eq:c2r-base-RT}, the corrected $RT$ coefficient now lies in
$I$.  All radial cross coefficients still contain an explicit factor $u$.

\subsection{Weighted boundedness and the strict residue}
Replace the radial root by the pure-coefficient rescaling
\[
T_N=\sqrt N\,(u^2-c^2).
\]
We check every type of coefficient of the resulting Gram.  The hard
block and the hard--$R$ block were already bounded.  A hard--$T$
coefficient contains $u/\sqrt{e_i}$, times bounded factors, for
$i=1,2,4,5$: this follows from the odd-parity normal division, the pure
current elimination, and the two displayed zero corrections.  Dividing
it by $\sqrt N$ remains bounded, with zero residue, because
\[
\frac1{\sqrt{Ne_1}}=\frac{\sqrt{r/N}}{\sqrt{U_1}},\quad
\frac1{\sqrt{Ne_2}}=\frac{\sqrt{g/N}}{\sqrt{U_2}},\quad
\frac1{\sqrt{Ne_4}}=\frac{\sqrt{h/N}}{\sqrt{U_4}},\quad
\frac1{\sqrt{Ne_5}}=\frac{\sqrt{j/N}}{\sqrt{U_5}}.
\]
The corrected $RT$ coefficient lies in $I$ and contains $u$, so its
quotient by $\sqrt N$ is bounded and has zero residue.

For the $T^2$ coefficient, the graph identity $f_*(u,u,1)=0$ says that
its monic scalar quotient $C_2(u^2)$ vanishes modulo $I$.  Every other
$T^2$ contribution contains $Q$ or $h+a^2r$, hence belongs to $I$ too.
Its quotient by $N$ is therefore bounded.  All such extra terms have
an explicit factor $u^2$.  At $u=0$ the scalar quotient is exactly
$4\alpha/c^4$, so \eqref{eq:c2r-weight-bounds} makes the weighted radial
pivot equal to $1$ at residue.  In particular the new correction adds
$2u^2(h+a^2r)/(D_cN)$ to that coefficient, which is bounded and has
zero residue.

The residue of the final Gram, in the roots
$(K_1,K_2,K_4,K_5,R,T_N)$, is therefore
\begin{equation}
4K_1^2+4K_2^2+4K_4^2+4K_5^2
 +(1+\overline W)^2R^2+T_N^2\succ0.
\label{eq:c2r-strict-residue}
\end{equation}
Here the relation $\zeta$ has residue $-(1+\overline W)R$, since all
four inverse packet scales are infinitesimal.  The strict-residue Gram
lemma proves the full bounded $Z=1,W\geq0$ chart.  The $W=1,Z$
infinitesimal proof of the preceding section also applies unchanged:
it uses boundedness, not a positive residue, of $e_3$.  Its four packet
roots have strict residue.  Sign dispatch and homogeneous restoration
therefore prove the theorem for all projective states.

\begin{corollary}[Complete bounded-shape C2 source family]
\label{cor:c2-bounded-shape}
In the unchanged C2 source let $c>0$ be infinitesimal, $r,g,j>0$
bounded, and $a,1-r,t$ positive bounded units.  Assume only
$h=b-a^2t>0$.  Then every independent zero-zero test and every projective
state have tensor-SOS parent.
\end{corollary}
\begin{proof}
The gap $h$ is bounded.  If it is a unit,
\eqref{eq:c2y-N-positive} makes $N$ a unit, so Theorem~\ref{thm:c2u}
applies.  If $h$ is infinitesimal and one of $r,g,j$ is a unit, use
Theorem~\ref{thm:c2-y-current}.  If all four are infinitesimal, use
Theorem~\ref{thm:c2-radial-corner}.  These alternatives are exhaustive.
\end{proof}
